\documentclass[10pt,a4paper]{amsart}
\usepackage{amsmath,amssymb,amsthm,graphicx}
% For commutative diagrams we use Xy-pic
\usepackage[all]{xy}

% We use 2cell for 2-commutative diagrams.
\xyoption{2cell}
% Theorem environments.
%
\theoremstyle{plain}
\newtheorem{theorem}[subsection]{Theorem}
\newtheorem{proposition}[subsection]{Proposition}
\newtheorem{lemma}[subsection]{Lemma}

\theoremstyle{definition}
\newtheorem{definition}[subsection]{Definition}
\newtheorem{example}[subsection]{Example}
\newtheorem{exercise}[subsection]{Exercise}
\newtheorem{situation}[subsection]{Situation}

\theoremstyle{remark}
\newtheorem{remark}[subsection]{Remark}
\newtheorem{remarks}[subsection]{Remarks}

\numberwithin{equation}{subsection}

% Macros
%
\def\lim{\mathop{\mathrm{lim}}\nolimits}
\def\colim{\mathop{\mathrm{colim}}\nolimits}
\def\Spec{\mathop{\mathrm{Spec}}}
\def\Hom{\mathop{\mathrm{Hom}}\nolimits}
\def\Ext{\mathop{\mathrm{Ext}}\nolimits}
\def\SheafHom{\mathop{\mathcal{H}\!\mathit{om}}\nolimits}
\def\SheafExt{\mathop{\mathcal{E}\!\mathit{xt}}\nolimits}
\def\Sch{\mathit{Sch}}
\def\Mor{\mathop{\mathrm{Mor}}\nolimits}
\def\Ob{\mathop{\mathrm{Ob}}\nolimits}
\def\Sh{\mathop{\mathit{Sh}}\nolimits}
\def\NL{\mathop{N\!L}\nolimits}
\def\CH{\mathop{\mathrm{CH}}\nolimits}
\def\proetale{{pro\text{-}\acute{e}tale}}
\def\etale{{\acute{e}tale}}
\def\QCoh{\mathit{QCoh}}
\def\Ker{\mathop{\mathrm{Ker}}}
\def\Im{\mathop{\mathrm{Im}}}
\def\Coker{\mathop{\mathrm{Coker}}}
\def\Coim{\mathop{\mathrm{Coim}}}

% Boxtimes
%
\DeclareMathSymbol{\boxtimes}{\mathbin}{AMSa}{"02}

%
% Macros for moduli stacks/spaces
%
\def\QCohstack{\mathcal{QC}\!\mathit{oh}}
\def\Cohstack{\mathcal{C}\!\mathit{oh}}
\def\Spacesstack{\mathcal{S}\!\mathit{paces}}
\def\Quotfunctor{\mathrm{Quot}}
\def\Hilbfunctor{\mathrm{Hilb}}
\def\Curvesstack{\mathcal{C}\!\mathit{urves}}
\def\Polarizedstack{\mathcal{P}\!\mathit{olarized}}
\def\Complexesstack{\mathcal{C}\!\mathit{omplexes}}
% \Pic is the operator that assigns to X its picard group, usage \Pic(X)
% \Picardstack_{X/B} denotes the Picard stack of X over B
% \Picardfunctor_{X/B} denotes the Picard functor of X over B
\def\Pic{\mathop{\mathrm{Pic}}\nolimits}
\def\Picardstack{\mathcal{P}\!\mathit{ic}}
\def\Picardfunctor{\mathrm{Pic}}
\def\Deformationcategory{\mathcal{D}\!\mathit{ef}}
\usepackage{cleveref}
\crefname{originaldefinition}{original definition}{original definitions}
\Crefname{originaldefinition}{Original definition}{Original definitions}
\makeatletter
\newlabel{definition-tau-site}{{03XB}{source}{}{sourcepointer-1}{}}
\newlabel{definition-tau-site@cref}{{[originaldefinition][0][]03XB}{[1][0][]source}{}{}{}}
\makeatother
\crefname{originallemma}{original lemma}{original lemmas}
\Crefname{originallemma}{Original lemma}{Original lemmas}
\makeatletter
\newlabel{topologies-lemma-contained-in}{{022J}{source}{}{sourcepointer-2}{}}
\newlabel{topologies-lemma-contained-in@cref}{{[originallemma][0][]022J}{[1][0][]source}{}{}{}}
\makeatother
\crefname{originallemma}{original lemma}{original lemmas}
\Crefname{originallemma}{Original lemma}{Original lemmas}
\makeatletter
\newlabel{sets-lemma-what-is-in-it}{{000R}{source}{}{sourcepointer-3}{}}
\newlabel{sets-lemma-what-is-in-it@cref}{{[originallemma][0][]000R}{[1][0][]source}{}{}{}}
\makeatother
\crefname{originallemma}{original lemma}{original lemmas}
\Crefname{originallemma}{Original lemma}{Original lemmas}
\makeatletter
\newlabel{topologies-lemma-alternative}{{04HR}{source}{}{sourcepointer-4}{}}
\newlabel{topologies-lemma-alternative@cref}{{[originallemma][0][]04HR}{[1][0][]source}{}{}{}}
\makeatother
\title{The etale topos of a scheme does not depend on the choice of universe for the small etale site}
\author{The Stacks Project Authors}
\date{stacks-project:etale-cohomology.tex:local-snapshot-20261011; GNU FDL 1.2 or later}
\begin{document}
\maketitle
\noindent\textit{Editorial note.} All delivered mathematical text is the author's own, copied from the pinned Stacks Project chapter file etale-cohomology.tex; no statement or proof step was rewritten, reordered or removed.\par\medskip
\noindent\textit{Editorial note.} The preamble is the chapter's own native preamble excerpt used for this source: the Xy-pic setup at lines 14-18 and the theorem declarations with the mathematical macros at lines 154-217 of sources/stacks-etale-cohomology/preamble.tex. The quoted class line and the xr-hyper externaldocument block are not included because the assembler supplies a plain amsart wrapper; no macro was invented or redefined.\par\medskip
\noindent\textit{Editorial note.} The same-chapter reference is located at its \textbackslash{}label line inside the pinned chapter file; the three cross-chapter references are located in the project's own tag list staged in the snapshot (sources/stacks-chow/tags.txt), at the verbatim lines carrying their canonical tags.\par\medskip
\noindent\textit{Allowed standard prerequisites (editorial).}\begin{itemize}
\item Construction of the small etale site of a scheme from a big etale site. Same chapter, Definition environment labelled definition-tau-site (source line 1909, Stacks tag 03XB); the statement refers to it explicitly. Supplies the construction whose independence from the chosen universe is asserted; the comparison argument is the counted content.
\item Reduction of a comparison of sites to a comparison of small categories and coverings. Topologies, Lemma topologies-lemma-contained-in (022J) and Lemma topologies-lemma-alternative (04HR), plus Sets, Lemma sets-lemma-what-is-in-it (000R); invoked at proof lines 2138, 2140 and 2149. Supply the reduction to the case of nested big sites, the smallness fact that every affine scheme etale over S sits in both, and the criterion that an equivalence preserving standard coverings induces an equivalence of topoi.
\end{itemize}\medskip
\ifdefined\hypertarget\hypertarget{sourcepointer-1}{}\fi
\noindent\textit{Source reference.} Original definition 03XB: sources/stacks-etale-cohomology/pinned-original-source.tex, source lines 1909--1909. Same-chapter pointer named in the counted statement (line 2131) for the construction whose independence is asserted; the comparison argument is retained verbatim in the counted span.\par\medskip
\ifdefined\hypertarget\hypertarget{sourcepointer-2}{}\fi
\noindent\textit{Source reference.} Original lemma 022J: sources/stacks-chow/tags.txt, source lines 2708--2708. Cross-chapter standard prerequisite invoked once at proof line 2138 to reduce to the case of nested big etale sites.\par\medskip
\ifdefined\hypertarget\hypertarget{sourcepointer-3}{}\fi
\noindent\textit{Source reference.} Original lemma 000R: sources/stacks-chow/tags.txt, source lines 36--36. Cross-chapter standard prerequisite invoked once at proof line 2140 for the smallness statement that an affine scheme etale over S belongs to both big sites.\par\medskip
\ifdefined\hypertarget\hypertarget{sourcepointer-4}{}\fi
\noindent\textit{Source reference.} Original lemma 04HR: sources/stacks-chow/tags.txt, source lines 5877--5877. Cross-chapter standard prerequisite invoked once at proof line 2149 for the criterion that an equivalence of small categories preserving standard coverings induces an equivalence of topoi; the verification of that preservation is retained in the counted span.\par\medskip
\begin{lemma}
\label{lemma-etale-topos-independent-partial-universe}
Let $S$ be a scheme. The \'etale topos of $S$ is independent
(up to canonical equivalence) of the construction of the small
\'etale site in Definition \ref{definition-tau-site}.
\end{lemma}

\begin{proof}
We have to show, given two big \'etale sites
$\Sch_\etale$ and $\Sch_\etale'$ containing
$S$, then $\Sh(S_\etale) \cong \Sh(S_\etale')$
with obvious notation. By Topologies, Lemma \ref{topologies-lemma-contained-in}
we may assume $\Sch_\etale \subset \Sch_\etale'$.
By Sets, Lemma \ref{sets-lemma-what-is-in-it}
any affine scheme \'etale over $S$ is isomorphic to an object
of both $\Sch_\etale$ and $\Sch_\etale'$.
Thus the induced functor
$S_{affine, \etale} \to S_{affine, \etale}'$
is an equivalence. Moreover, it is clear that both this functor
and a quasi-inverse map transform standard \'etale coverings into
standard \'etale coverings.
Hence the result follows from
Topologies, Lemma \ref{topologies-lemma-alternative}.
\end{proof}

\end{document}
